\section*{Hoofdstuk \ref{ch:b3:gravity-geodesy} --- Zwaartekracht en de vorm van de Aarde}

\begin{solution}{exo:b3:gravity-geodesy:1}
$b = a(1 - f) = 6378137\times0.99664719 = \qty{6356752.3}{m}$, en
$a - b = \qty{21384.7}{m}$, ongeveer \qty{21.4}{km}. Clairaut:
$f = (3J_2 + m)/2 = (0.00324789 + 0.00344979)/2 = \num{0.0033488}$, of
$1/298.61$, tegen $1/298.26$: \qty{0.12}{\percent} kleiner, de grootte van de
termen van tweede orde die de relatie verwaarloost.
\end{solution}

\begin{solution}{exo:b3:gravity-geodesy:2}
$\sigma = \qty{6.0e13}{kg}/\qty{5.0e11}{m^2} = \qty{120}{kg.m^{-2}}$, dus
$h_w = 120/1000 = \qty{0.12}{m}$: \qty{12}{cm} water. Verdeeld over de
oceaan, $60/360 = \qty{0.17}{mm}$ van het zeeniveau.
\end{solution}

\begin{solution}{exo:b3:gravity-geodesy:3}
Aan de evenaar $\gamma = \qty{9.780327}{m.s^{-2}}$. Bij \ang{45},
$\sin^2\varphi = 0.5$ en $\sin^22\varphi = 1$: $\gamma = 9.780327\times1.0026454
= \qty{9.806200}{m.s^{-2}}$. Aan de polen, $9.780327\times1.0053024 =
\qty{9.832186}{m.s^{-2}}$. Het verschil is \qty{0.051859}{m.s^{-2}}, of
\qty{5186}{\milli\gal}.
\end{solution}

\begin{solution}{exo:b3:gravity-geodesy:4}
$D = ET_e^3/(12(1 - \nu^2)) = 10^{11}T_e^3/11.25$: \qty{8.89e21}{N.m} voor
\qty{10}{km}, \qty{2.40e23}{N.m} voor \qty{30}{km}. Met $(\rho_m - \rho_w)g =
\qty{22269}{N.m^{-3}}$, $\alpha = (4D/22269)^{1/4}$: \qty{35.5}{km} en
\qty{81.0}{km}. De greppel eindigt bij $3\pi\alpha/4$, \qty{84}{km} en \qty{191}{km};
de toppen van de \omterm{def:b3:gravity-geodesy:alpha}{voorwelving} liggen bij $\pi\alpha$, \qty{112}{km} en \qty{255}{km}.
\end{solution}

\begin{solution}{exo:b3:gravity-geodesy:5}
(a) $T = -\gamma N = -9.80\times50 = \qty{-490}{J.kg^{-1}}$. (b) De geoïde kantelt
met $10/100\,000 = \num{1.0e-4}$ rad ten opzichte van de ellipsoïde, en de schietlijn,
loodrecht op de geoïde, met dezelfde hoek ten opzichte van de normaal op de ellipsoïde:
$\num{1.0e-4}\times206\,265 = 20.6''$.
\end{solution}

\begin{solution}{exo:b3:gravity-geodesy:6}
$H = J_2/(C/Ma^2) = \num{1.08263e-3}/0.3308 = \num{0.0032728}$, en
$1/H = 305.55$. $J_2$ is per hypothese hetzelfde; een homogene Aarde zou
$C/(Ma^2) = 0.4$ hebben en dus $H = J_2/0.4 = \num{0.0027066}$: de precessie zou
trager verlopen. Door zowel $J_2$ als $H$ te meten wordt de concentratie van massa
naar het centrum toe onthuld.
\end{solution}

\begin{solution}{exo:b3:gravity-geodesy:7}
$r = \qty{6878.137}{km}$; $n = \sqrt{GM/r^3} = \qty{1.10678e-3}{rad.s^{-1}}$, een
periode van $2\pi/n = \qty{94.6}{min}$. $\dot\Omega = -1.5\times\num{1.10678e-3}
\times\num{1.08263e-3}\times(6378.137/6878.137)^2\times\cos\ang{89} =
\qty{-2.70e-8}{rad.s^{-1}}$, of \ang{-0.134} per dag, westwaarts. Een bijna-polaire baan
passeert elke breedtegraad, zodat de hele Aarde in kaart wordt gebracht; ze precesseert ook
langzaam, zodat de grondsporen zich regelmatig aanvullen.
\end{solution}

\begin{solution}{exo:b3:gravity-geodesy:8}
$e^{-2\pi\times500/100} = e^{-31.4} = \num{2.3e-14}$: \qty{2e-13}{\milli\gal},
verwaarloosbaar. $e^{-2\pi\times500/1000} = e^{-3.14} = 0.0432$:
\qty{0.43}{\milli\gal}.
\end{solution}

\begin{solution}{exo:b3:gravity-geodesy:9}
$3\pi\alpha/4 = \qty{120}{km}$ geeft $\alpha = \qty{50.9}{km}$.
$D = (\rho_m - \rho_w)g\alpha^4/4 = 22269\times(50\,930)^4/4 =
\qty{3.75e22}{N.m}$, en $T_e = (11.25D/E)^{1/3} = \qty{16.2}{km}$. De plaat
kruist zijn ongebogen niveau bij $\pi\alpha/2 = \qty{80}{km}$ vanaf de trog.
\end{solution}

\begin{solution}{exo:b3:gravity-geodesy:10}
(a) Met $u = x/\alpha$, $\mathrm d^4/\mathrm dx^4 = \alpha^{-4}\,\mathrm
d^4/\mathrm du^4$. Door $\cos u + \sin u = \sqrt2\cos(u - \pi/4)$ te schrijven, is de
functie $\sqrt2w_0\,\mathrm{Re}\,e^{(-1 + i)u - i\pi/4}$, en
$(-1 + i)^4 = -4$: dus $w'''' = -4w/\alpha^4$. Dan $Dw'''' + (\rho_m -
\rho_i)gw = (-4D/\alpha^4 + (\rho_m - \rho_i)g)w = 0$, omdat $\alpha^4 =
4D/((\rho_m - \rho_i)g)$. (b) $w'' = -(2w_0/\alpha^2)e^{-u}(\cos u - \sin u)$, dus
$M = -Dw'' = (2Dw_0/\alpha^2)e^{-u}(\cos u - \sin u)$: grootst in absolute waarde bij
$u = 0$, nul waar $\tan u = 1$, $x = \pi\alpha/4$. (c) Onder de belasting is de
plaat naar boven hol, haar bovenvlak onder compressie; voorbij $\pi\alpha/4$
is zij naar boven bol, haar bovenvlak uitgerekt. Daar $M' \propto
-e^{-u}\cos u$ verdwijnt bij $u = \pi/2$: de uitrekking is het grootst bij
$x = \pi\alpha/2$, aan de buitenflank van de greppel richting de \omterm{def:b3:gravity-geodesy:alpha}{voorwelving}, waar
$|M|$ bedraagt $e^{-\pi/2} = \qty{21}{\percent}$ van zijn waarde onder de belasting. Dat is
waar afschuivingen en scheuren zich aan de bovenkant van de plaat zouden openen.
\end{solution}

\begin{solution}{exo:b3:gravity-geodesy:11}
$\sqrt{1 + \eta} = 2.5(1 - 1.5\times0.3308) = 1.2595$, dus $1 + \eta = 1.5863$
en $\tfrac52m/f = 2.5863$: $f_h = 2.5\times0.00344979/2.5863 = \num{0.0033346}$,
$1/f_h = 299.9$. Het verschil $f - f_h = \num{1.82e-5}$; vermenigvuldigd met $a$ ligt de
evenaar \qty{116}{m} boven de hydrostatische figuur. De Aarde is niet helemaal
een vloeistof in evenwicht: het inwendige handhaaft spanningen, afkomstig van de sterkte van
de lithosfeer of van de stromings- en dichtheidsverschillen in de mantel, die
de vorm een beetje van het evenwicht verwijderd houden.
\end{solution}

\begin{solution}{exo:b3:gravity-geodesy:12}
De knoop moet vooruitgaan met $2\pi/(365.242\times86\,400) =
\qty{1.991e-7}{rad.s^{-1}}$. Bij $r = \qty{7078.137}{km}$, $n =
\qty{1.0602e-3}{rad.s^{-1}}$ en $\tfrac32nJ_2(a/r)^2 = \num{1.398e-6}$, dus
$\cos i = -\num{1.991e-7}/\num{1.398e-6} = -0.1424$ en $i = \ang{98.2}$: een
licht retrograde, zonsynchrone baan.
\end{solution}

\begin{solution}{pb:b3:gravity-geodesy:1}
\textbf{1.} $m = (\num{7.292115e-5})^2\times6378137^2\times6356752.3/
\num{3.986005e14} = \num{0.00344979}$.

\textbf{2.} $f = (3J_2 + m)/2 = \num{0.0033488}$, $1/f = 298.61$.

\textbf{3.} $\sqrt{1 + \eta} = 2.5(1 - 1.5\times0.3308) = 1.2595$,
$\tfrac52m/f_h = 2 + 0.5863$, dus $f_h = \num{0.0033346}$ en $1/f_h = 299.88$.

\textbf{4.} De Clairaut-waarde, $1/298.61$, ligt binnen \qty{0.12}{\percent} van
de waargenomen $1/298.26$; de hydrostatische, $1/299.88$, is kleiner met
$\Delta f = \num{1.82e-5}$, en $a\,\Delta f = \qty{116}{m}$: de evenaar ligt
\qty{116}{m} boven de hydrostatische figuur.

\textbf{5.} $H = \num{1.08263e-3}/0.3308 = \num{0.0032728}$ ($1/305.6$).

\textbf{6.} $\tfrac12\times\qty{80}{km}\times\qty{6}{km} = \qty{240}{km^2} =
\qty{2.40e8}{m^2}$.

\textbf{7.} $V_0 = (2800 - 1030)\times9.81\times\num{2.40e8} =
\qty{4.17e12}{N.m^{-1}}$.

\textbf{8.} Airy: $(\rho_m - \rho_v)r = (\rho_v - \rho_w)h$, dus $r =
6\times1770/500 = \qty{21.2}{km}$ onder de kam.

\textbf{9.} De zeebodem wordt ver voorbij de voet van de rug doorgebogen: de
belasting wordt niet kolom voor kolom gecompenseerd, maar door de buiging van de plaat over een
gebied dat viermaal zo breed is. De plaat heeft sterkte: regionale compensatie.

\textbf{10.} $\alpha = 165/(3\pi/4) = \qty{70.0}{km}$.

\textbf{11.} $\alpha = 220/\pi = \qty{70.0}{km}$: hetzelfde, zoals een enkele elastische
plaat vereist.

\textbf{12.} $D = (\rho_m - \rho_w)g\alpha^4/4 = 2270\times9.81\times(\num{7.00e4})^4/4
= \qty{1.34e23}{N.m}$.

\textbf{13.} $T_e = (12(1 - \nu^2)D/E)^{1/3} = (11.25\times\num{1.34e23}/
\num{1e11})^{1/3} = \qty{24.7}{km}$.

\textbf{14.} $w_0 = V_0\alpha^3/(8D) = \num{4.17e12}\times(\num{7.00e4})^3/
(8\times\num{1.34e23}) = \qty{1334}{m}$.

\textbf{15.} $w_0e^{-\pi} = 1334\times0.0432 = \qty{58}{m}$.

\textbf{16.} $2\times1334\times\num{7.00e4}\times2270\times9.81 =
\qty{4.16e12}{N.m^{-1}}$, gelijk aan $V_0$ binnen de afrondingsfout van $D$.

\textbf{17.} \qty{24.7}{km} ligt binnen het bereik van de Hawaïaanse Rug. De
plaat zakt \qty{1.3}{km} onder de as in plaats van een wortel van
\qty{21}{km} te vormen: dezelfde opwaartse kracht is verdeeld over de \qty{330}{km} van de geul.

\textbf{18.} $\lambda = 2\pi\times\qty{70.0}{km} = \qty{440}{km}$,
$k = 1/\alpha = \qty{1.43e-5}{m^{-1}}$.

\textbf{19.} $\sigma = (2900 - 1030)\times500 = \qty{9.35e5}{kg.m^{-2}}$, dus
$2\pi G\sigma = \qty{3.92e-4}{m.s^{-2}} = \qty{39.2}{\milli\gal}$ op de zeebodem;
vermenigvuldigd met $e^{-k\times\qty{5}{km}} = 0.931$, \qty{36.5}{\milli\gal} op zeeniveau
(negatief waar de plaat naar beneden wordt gedrukt).

\textbf{20.} $\sigma = (3300 - 2900)\times500 = \qty{2.0e5}{kg.m^{-2}}$:
\qty{8.39}{\milli\gal}, vermenigvuldigd met $e^{-k\times\qty{12}{km}} = 0.843$,
\qty{7.07}{\milli\gal}. Het totaal is \qty{43.6}{\milli\gal}.

\textbf{21.} $e^{-k\times\qty{400}{km}} = e^{-5.71} = 0.0033$: \qty{0.14}{\milli\gal}.

\textbf{22.} $e^{-2\pi\times400/80} = e^{-31.4} = \num{2.3e-14}$: de satelliet
ziet niets van de rug zelf, alleen de brede buiging, en ook die is
driehonderdvoud verzwakt; de eigen anomalie van de rug moet op zee worden gemeten.

\textbf{23.} $n \approx 2\pi\times6371/440 = 91$.

\textbf{24.} Uit de vragen 10 tot en met 13 volgt dat de \omterm{def:b3:gravity-geodesy:flexure}{effectieve elastische dikte} van de
plaat onder de Oru-rug
\textbf{\qty[text-series-to-math=true]{24.7}{km}} is.
\end{solution}
